\section{Kerouac and the Faustian West}

\begin{quotex}
The sinner is at the very heart of Christianity. Nobody is so competent as the sinner in matters of Christianity. Nobody, except the saint. \flright{\textsc{Charles Péguy}}

Love, Work, and Suffer \flright{Motto of the Kerouac family (\emph{Rivista Araldica})}

Am actually not ``beat" but strange solitary crazy Catholic mystic \flright{\textsc{Jack Kerouak}, \emph{Lonesome Traveler}}

\end{quotex}
Since I was under the impression that Jack Kerouac was lost in the memory hole, I was recently surprised by the publication of \emph{Jack Kerouac and the Decline of the West}\footnote{\url{https://shop.aer.io/roguescholar}}, an essay by \textbf{Semmelweis}, published by Rhodes Scholar Press.

Since ``Semmelweis" claims to be a GenX latchkey kid, the lifestyle described by Kerouac is hardly a living option for him. For my generation (I have GenX sons), Kerouac, when was even acknowledged, was either a passing ``stage" one passed through, a proto-hippy, or lumped in, rather inaccurately, with the ``beats". Semmelweis, unburdened by such preconceptions, is able to separate the ``real" Kerouac from the stereotypes of the beatnik … not that that task is so easy to do, since Kerouac was rather complex and his behavior did not always match his innermost thoughts. What Semmelweis is able to see is much deeper than what those of Kerouac's generation were able to see.

\paragraph{One Man and Three Respectable People}
\begin{quotex}
Favorite complaint about contemporary world: the facetiousness of ``respectable" people \flright{\emph{Lonesome Traveler}}

\end{quotex}
Semmelweis comments on Kerouac's appearance on \emph{Firing Line}\footnote{\url{https://www.youtube.com/watch?v=BYgv7ur8ipg}} in 1968, which is worth watching.

Of the four men on stage, Kerouac was the only one who was able to connect with the audience. Allegedly intoxicated, he brought the audience to laughter several times, while the three ``respectable men" droned on. They are what the French call a ``type"; that is, they are devoid of individual personality. These are the other three on the show:

\begin{itemize}
\item \textbf{William F. Buckley}: At that time, he was the Maxwell's Demon of conservatism, deciding who was in and who was out. After 50 years, his brand of conservatism has conserved nothing. 
\item \textbf{Lewis Yablonsky}: Professor Yablonsky, I should add. He was the ``expert" on hippies, although he contributed nothing. There was no creed for hippies nor membership cards, so he contributed nothing. 
\item \textbf{Ed Sanders}: He was the ``hippie" type, and the contrast with Kerouac could not have been clearer. He recited the platitudes of love and peace, etc., like a cardboard cutout. You can see them in the baseball games this week. Put the cutout on a new show in 2020 and you would not even know the difference. The same speech repeated ad nauseum, and everyone thinks it is new.

In high school, I had the Fugs album. When my mother heard it, she took it off the record player and returned it to the record store for an exchange. I guess she could see further than I could. 
\end{itemize}
\paragraph{Personal Encounter}
A generation ahead of me, Kerouac grew up in Lowell, Massachusetts, not very far from my home town. As a rather bookish youth, I tended to act out roles from the books I read. Not quite willing and able to takes things as far as Kerouac did, I did find opportunities. So in some ways, I understand his mindset from the inside, not just descriptively. Fortunately for me, I don't have an addictive type personality and was left unscarred. Unfortunately, that encounter has left me with an attraction to troubled women, like the sad but beautiful Tristessa.

My boomer friends, feeling the cold breath of impending death on their necks, have grown nostalgic, as though memories of old times will be as restorative as blood plasma transfusions from youths. So they send me links to albums by an 80 year old Dion, or old tracks from the wrinkled Rolling Stones. I'd be much more impressed if Dion became an anchorite or Jagger went full sannyasi.

Few of them have had a new idea since they were 19, despite professional successes. In their best moments, they sound like a second rate Ed Sanders. They can't figure me out, but I have just followed the logic.

There is a story, perhaps apocryphal, that some hippies tracked him down in St. Petersburg and went to visit to pay homage. He threw them out while claiming to be Catholic and conservative. The hippies didn't get the point, although some of us did.

\paragraph{Man as adventurer}
\begin{quotex}
Kerouac never lost his longing for a loving family, for a wife and children, and for the simpler, wholesome life of small town America. But he was always thwarted by his longing for adventure, passion, excitement, and lust. \flright{\textsc{Semmelweis}}

\end{quotex}
If you watch some youtube videos today, you may come to the conclusion that to be religious, it is only necessary to wear a scapula, pray so many rosaries, etc. This is not a criticism of those men or their practices, especially since in this day and age those are acts of defiance. Yet through Christian history, there has always been the desire for more. There have been explorers, knights, traders, warriors, pilgrims, missionaries, all with the urge for travel and adventure. They kept the faith, no matter how imperfectly.

Most of those options are not available today, so Kerouac made America his exploration (and more followed later). A GenXer cannot imagine what the USA was like in 1956, either physically or psychologically. That description will have to wait for my autobiography.

It seems that converts and reverts have been dominating the public discussion of religion. They are sincere and enthusiastic, and know their dogmas, canon law, and rituals perfectly. However, they often lack a certain ``feel", a Catholic mind that has embraced centuries, the globe, and been encultured with tales of saints, sinners, mystics, philosophers, and so on. There are classical pianists who are technically proficient and know how to play each note at the right time. Nevertheless, they lack an aesthetic ``feel" that gives life to the music, so they never make it to the top ranks.

Kerouac, on the other hand, has that ``feel", even if he is not a good role model for your children. For example, while getting high with William Burroughs in Morocco, he could be inspired by a Muslim poem, or image some old man on beach as the coming Bodhisattva. Yet in his best moments, he could have a genuine spiritual experience.

\begin{quotex}
And on Good Friday afternoon a heavenly performance of the St Matthew's Passion … I cried most of the time and had a vision of an angel in my mother's kitchen … and I realized it didn't matter that we sin, that all my own petty gripes didn't matter either. \flright{\emph{Lonesome Traveler}}

\end{quotex}
\paragraph{Final Plans}
\begin{quotex}
hermitage in the woods, quiet writing of old age, mellow hopes of Paradise \flright{\emph{Lonesome Traveler}}

\end{quotex}
Kerouac never got to be old, since alcohol killed him. That should be the goal of every old man, and Kerouac knew it. When Siddhartha tired of the pursuit of money, women, family, he, too, isolated himself. But he could think, fast, and wait. Jack was addicted and impatient, so his road of excess cannot be recommended, even if it leads to the palace of wisdom. Did he attain Paradise? Perhaps, since Paradise is for saints and sinners, not the lukewarm.

\paragraph{The peasant farmers}
Semmelweis picks out an interesting theme in Kerouac that I never would have noticed. He references the Fellaheen or ``peasant farmer": those who persist after a civilization has collapsed. Oswald Spengler defines them this way.

\begin{quotex}
only the primitive blood remains, alive, but robbed of its strongest and most promising elements. This residue is the Fellah type … Life as experienced by primitive and by fellaheen peoples is just the zoological up and down, a planless happening without goal or cadenced march in time, wherein occurrences are many, but, in the last analysis, devoid of significance." \flright{\textsc{Oswald Spengler}, \emph{The Decline of the West}}

\end{quotex}
Valentin Tomberg attributes them to the forgetting of the past. In our time, we can see that the forgetting is deliberate.

\begin{quotex}
the ``primitive" tribes and nomadic peoples, disinherited from their past and obliged to begin everything again, began to live in caves or camp under trees. There were once powerful kingdoms and magnificent towns but their descendants had lost all memory of them and gave themselves up entirely to the daily life of ``primitive" tribes —the life of hunting, fishing, agriculture and war. \flright{\textsc{Valentin Tomberg}, \emph{Meditations on the Tarot}}

\end{quotex}
Semmelweis points out that Kerouac identified with the Fellaheen of his time: hustlers, drug addicts, winos, even though they could not sustain even a primitive civilization. Nevertheless, Semmelweis points out the logic: Beat = beatific, therefore the Fellaheen, the lumpen, the outcasts, are holy. Perhaps this might also be due to the way the respectable men of the publishing industry marginalized him, so he was never able to penetrate into the respectable literary circles.

\paragraph{Faustian and Magian Civilizations}
Besides the Fellaheen, Semmelweis pulls out another point from Spengler. Kerouac was prophetic about the 21st century. Semmelweis explains:

\begin{quotex}
He saw the distinction in Spenglerian terms, which classifies Western European civilization as ``Faustian," and Near-Eastern civilization as ``Magian." 

\end{quotex}
It's remarkable that he saw that 21st century culture and spirituality would not be ``American,"

\begin{quotex}
which in this context means Western European and Faustian, but would be ``Magian," of the East, and of a type which rather than promoting the heroism and individualism of Faustian man, promotes the dissolution of the individual ego into the greater collective Spirit. 

\end{quotex}
That is why the spirituality of the East involves the dissolution of the Person in the unconditioned state, while Western spirituality, like in Catholicism, there is always the two in one. That so many Westerners are preferring Eastern spirituality is just a regrettable sign of the times.

\paragraph{Dionysus vs Dionysius}
\begin{quotex}
I believe in order, tenderness, and piety. 

\end{quotex}
Kerouac confessed this to Buckley:

\begin{quotex}
[the counterculture is] apparently some kind of Dionysian movement in late civilization, and which I did not intend, any more than I suppose Dionysus did … although I'm not Dionysius the Areopagite. I should have been. 

\end{quotex}
Semmelweis expands on that distinction:

\begin{quotex}
The point [is] that Kerouac's sensibilities and values are religious like Dionysius the Areopagite, not chaotic and destructive like the god Dionysus. 

\end{quotex}
Since the two names differ by only one iota, it shows how close the two temptations are. At birth, we are assigned a good angel and a bad angel. Hence, life is a perpetual spiritual warfare; that is not an option and you cannot be a draft dodger.

\paragraph{Errata}
The beatnik on the TV Show, \emph{The Many Loves of Dobie Gillis}, was named Maynard G. Krebs.



\flrightit{Posted on 2020-07-27 by Cologero }

\begin{center}* * *\end{center}

\begin{footnotesize}\begin{sffamily}



\texttt{Max on 2020-07-27 at 15:04 said: }

In Jünger's novel Eumeswil, the only non-fellaheen part of the population left seems to be the tyrant and his entourage. It describes a post-democratic polity, in the Platonic sense. Now old and tired, having already cycled through civilization, those who still possesses some measure of vitality reluctantly governs the rest. You do not have to agree in everything, but their company is still proves more interesting and valuable than the small-minded people down in the city. Of course, towards the end, those more energetic elements, however flawed they might be, will leave on a mysterious expedition, never to be seen again. Religious sensibility without a religion. Apollonian without an Apollo, but still better off than the converse. The fellaheen probably hardly notices one way or the other, and could not care less. There is no point in identifying with those left behind, or when the time comes, to look back.


\hfill

\texttt{Michael M on 2020-07-28 at 17:27 said: }

``Fortunately for me, I don't have an addictive type personality and was left unscarred. Unfortunately, that encounter has left me with an attraction to troubled women, like the sad but beautiful Tristessa."

Hey I thought you said you didn't have an addictive type personality! Addictive types of personalities seem to be called to the Spirit, what a powerful drive and a great weakness and impurity and what we don't realize in them: ``For this reason, as we shall afterwards say, God leads into the dark night those whom He desires to purify from all these imperfections so that He may bring them farther onward." Perhaps there's something to the ``addictive" element that no longer has a compass in our times, barring those with biological issues.

``We have a huge barrel of wine, but no cups.

That's fine with us. Every morning

we glow and in the evening we glow again.

They say there's no future for us. They're right.

Which is fine with us."

Beatnik or Sufi? Let the inner element of the reader decide, moreover let the inner element be more than passively witnessing.

Thank you for the article, the Saint and the Sinner, there's no room for error, an example of the trauma of the world and the light that penetrates through all, only then to transcend, though the closed circle continues, not to be celebrated but overcome.

```Cold, hunger, hate, mockery, scorn, injustice, prison, illness and even death?'

— I know it.'

— ‘Do you expect to be shunned by everyone? Do you expect to be totally alone?'

— ‘I am ready. I know it. I shall bear all the suffering and all the blows'.

— ‘Even if they do not come from enemies, but from parents, from friends?'

— ‘Yes… even from those…'

— ‘Good. Do you accept the sacrifice?'

— ‘Yes'.

— ‘An anonymous sacrifice? You will perish and nobody… but nobody will even know whose memory to honour?'

— ‘I have no use for recognition and pity. I have no use for a name.'

— ‘Are you ready for crime?' The young girl bowed her head.

— ‘Even for crime.

— ‘Do you know that one day you will believe no more in what you believe in now, and come to think that you have been a dupe and that it was for nothing that you have lost your young life?'

— ‘That too I know. Well though I know it, I wish to enter.' The young girl crossed the threshold. A heavy curtain fell. Gritting his teeth, someone uttered behind her:

— ‘A foolish girl!'

At which, from another place, a voice replied:"

I'll keep the cliffhanger this time.


\hfill

\texttt{AghorNath on 2020-07-28 at 19:26 said: }

Quite curious the timing of this post Cologero. 

Thank you for this, brings back memories of my days as a NYC youth in the early to mid-nineties; though I'm not looking for any restorative blood plasma transfusions on my stroll down memory lane. A GenXer myself, I developed a spiritual friendship in my late teens with someone of your generation who happened to be quite close to Kerouac at one time. Just last week I was informed of his death, though it's been several years now since his passing. He mentioned Kerouac and Ginsberg on a few occasions, just in passing. When we first became acquainted I thought he might be blowing smoke up my arse, all the people he would mention so casually, it was quite dizzying. That was until he introduced me to Peter Lamborn Wilson (who was one of Seyyed Hossein Nasr's disciples) when I had a pressing question about one of the Shi'ite Imams. I was inquiring about Islam at the time. Hassan (said friend), was himself Muslim, and apparently many years ago (long before I ever met him) a shaykh in a sufi order. Even though by the time we met, he had been battling addiction, there was always an air of saintliness permeating, irregardless of his vice. The expression ``give you the shirt off his back" doesn't do the man justice, at all. A sinner and a saint if there ever was one. Some would see just a washed up hipster, others an eccentric downtown grandfatherly figure. But I knew better, because while most knew only the exterior…drugs and Downtown NY life…I was probably one of the few who knew his True Man. And the regret he felt for living a life on the margins. He was wholly conscious of his sins. But he never forgot his Lord. 

He told me once, Kerouac and him had spent literally days locked in a tiny room discussing the nature of Christ (no drugs involved). And that the coffee shop crowd in alphabet city back in the day could ``never figure us out…we left them always scratching there heads…were we left-wing, right-wing(?)…labels are for the birds". Though he did always say him and his mates (Kerouac etc) where ``bohemian…NOT beats or hippies!"

After an old acquaintance told me of Hassan's passing, I googled his Christian name to see if there were any pictures floating about so I could say a prayer for him while in the presence of his image. A touching, if impersonal and lacking (on his spiritual life) article popped up that was written about him by The Daily Beast… 

\url{https://www.thedailybeast.com/from-the-beat-generation-to-c-squatfriends-bid-farewell-to-a-counterculture-forrest-gump}

It turns out that Kerouac put him in a book. Under the name ``Joey Rosenberg" in Big Sur. I heard tons of stories but this was news to me! Never read Kerouac, now is as good a time as any.

It seems that Mr. Kerouac and Mr. Heiserman (Hassan) may have suffered as the Holy Fool and Malamati respectively. Of the variety that there nature is kept hidden (to most)…even unto themselves. I always wondered (during our conversations, and even years later when i'd recall our talks) if he was a Bodhisattva, or my own personal avatara secretly guiding me out of troubled times. Probably not, but his deeply altruistic heart and his often sagacious advice in my youth gives me pause for thought.

Did he attain Paradise? The sinner is closer to God.

Inna lillahi wa inna ilayhi raji'un

P.s. – Praying your health is better Cologero… Godspeed


\hfill

\texttt{AghorNath on 2020-07-28 at 19:56 said: }

Also interestingly, I forgot to mention in my previous comment. Last night I had been watching a nice conversation of thinkers on spirituality and community. Speakers included Seyyed Hossein Nar and William F. Buckley. Buckley in his usual dry manner was something of an odd duck within the group not quite understanding why he was there and holding fast and firm to his exclusivist outlook. When other panel members would seemingly like to scoff at Buckley for his exclusivity, Nasr came to his defense asserting ``…No religion can accept the idea that it is not True. I agree 100\% with Mr. Buckley, if he rejects me as being a heretic, I nevertheless except his position, because he claims the Truth of Christianity is THE Truth, that's why he follows it. If we just have truth as relativity, then why do we have to attach ourselves to it. There is a sense of absoluteness in all religions…"

Anyway worth a watch \url{https://www.youtube.com/watch?v=aq7BJLPfiYw}


\hfill

\texttt{Cologero on 2020-07-30 at 17:22 said: }

That brings to mind the ``double truth" theory\footnote{\url{https://en.wikipedia.org/wiki/Double_truth}}, that philosophy and religion could be contradictory, yet both be true.


\hfill


\end{sffamily}\end{footnotesize}
